\documentclass[12pt, a4paper]{article}

\usepackage[utf8]{inputenc}
\usepackage[T1]{fontenc}
\usepackage[brazil]{babel}
\usepackage{amsmath, amssymb}
\usepackage{geometry}
\usepackage{booktabs}

\geometry{margin=2.5cm}

\title{\textbf{Propriedades de Somatórios}}
\author{Lucas Rafael Pessoa do Nascimento}
\date{}

\begin{document}

\maketitle

\section{Definição}

O somatório é uma notação compacta para representar a soma de uma sequência de termos:

\[
\sum_{i=m}^{n} f(i) = f(m) + f(m+1) + f(m+2) + \cdots + f(n)
\]

onde $i$ é o índice de somatório, $m$ é o limite inferior e $n$ é o limite superior.

% -------------------------------------------------------
\section{Propriedades Fundamentais}

\subsection{Linearidade}

\subsubsection*{Fator constante}
Uma constante multiplicativa pode ser retirada do somatório:

\[
\sum_{i=m}^{n} c \cdot f(i) = c \sum_{i=m}^{n} f(i)
\]

\subsubsection*{Aditividade}
A soma de duas funções pode ser separada em dois somatórios:

\[
\sum_{i=m}^{n} \bigl[f(i) + g(i)\bigr] = \sum_{i=m}^{n} f(i) + \sum_{i=m}^{n} g(i)
\]

\subsection{Somatório de Constante}
Quando o termo não depende do índice:

\[
\sum_{i=m}^{n} c = c \cdot (n - m + 1)
\]

\subsection{Partição do Intervalo}
Um somatório pode ser dividido em dois somatórios parciais, desde que $m \leq k < n$:

\[
\sum_{i=m}^{n} f(i) = \sum_{i=m}^{k} f(i) + \sum_{i=k+1}^{n} f(i)
\]

\subsection{Mudança de Índice}
O índice pode ser deslocado sem alterar o valor:

\[
\sum_{i=m}^{n} f(i) = \sum_{j=m+c}^{n+c} f(j - c)
\]

\subsection{Inversão da Ordem}
A ordem de somatório pode ser invertida:

\[
\sum_{i=m}^{n} f(i) = \sum_{i=m}^{n} f(m + n - i)
\]

% -------------------------------------------------------
\section{Fórmulas Fechadas Importantes}

\subsection{Soma dos $n$ primeiros inteiros positivos}

\[
\sum_{i=1}^{n} i = \frac{n(n+1)}{2}
\]

\textbf{Exemplo:} $\displaystyle\sum_{i=1}^{5} i = 1 + 2 + 3 + 4 + 5 = 15 = \frac{5 \cdot 6}{2}$

\subsection{Soma dos quadrados}

\[
\sum_{i=1}^{n} i^2 = \frac{n(n+1)(2n+1)}{6}
\]

\subsection{Soma dos cubos}

\[
\sum_{i=1}^{n} i^3 = \left[\frac{n(n+1)}{2}\right]^2
\]

\subsection{Progressão Geométrica (série geométrica finita)}
Para $r \neq 1$:

\[
\sum_{i=0}^{n} r^i = \frac{r^{n+1} - 1}{r - 1}
\]

Caso especial com $r = 2$:

\[
\sum_{i=0}^{n} 2^i = 2^{n+1} - 1
\]

\subsection{Série Harmônica (aproximação)}

\[
\sum_{i=1}^{n} \frac{1}{i} \approx \ln n + \gamma
\]

onde $\gamma \approx 0{,}5772$ é a constante de Euler-Mascheroni.

% -------------------------------------------------------
\section{Somatórios Duplos}

Um somatório duplo percorre duas variáveis independentes:

\[
\sum_{i=1}^{m} \sum_{j=1}^{n} f(i, j)
\]

\subsubsection*{Intercambialidade (limites independentes)}
Quando os limites não dependem um do outro, a ordem pode ser trocada:

\[
\sum_{i=1}^{m} \sum_{j=1}^{n} f(i, j) = \sum_{j=1}^{n} \sum_{i=1}^{m} f(i, j)
\]

% -------------------------------------------------------
\section{Aplicação em Análise de Algoritmos}

Somatórios são essenciais para calcular a complexidade de algoritmos. Exemplos comuns:

\begin{center}
\begin{tabular}{lll}
\toprule
\textbf{Estrutura} & \textbf{Somatório} & \textbf{Complexidade} \\
\midrule
Loop simples de $1$ a $n$ & $\displaystyle\sum_{i=1}^{n} 1$ & $\Theta(n)$ \\[10pt]
Loop com trabalho $i$ & $\displaystyle\sum_{i=1}^{n} i$ & $\Theta(n^2)$ \\[10pt]
Loop dividindo por 2 & $\displaystyle\sum_{i=0}^{\log n} 1$ & $\Theta(\log n)$ \\[10pt]
Dois loops aninhados & $\displaystyle\sum_{i=1}^{n}\sum_{j=1}^{i} 1$ & $\Theta(n^2)$ \\
\bottomrule
\end{tabular}
\end{center}

% -------------------------------------------------------
\section{Resumo das Propriedades}

\[
\sum_{i=m}^{n} c \cdot f(i) = c \sum_{i=m}^{n} f(i)
\qquad
\sum_{i=m}^{n}[f(i)\pm g(i)] = \sum_{i=m}^{n}f(i) \pm \sum_{i=m}^{n}g(i)
\]

\[
\sum_{i=1}^{n} i = \frac{n(n+1)}{2}
\qquad
\sum_{i=1}^{n} i^2 = \frac{n(n+1)(2n+1)}{6}
\qquad
\sum_{i=0}^{n} r^i = \frac{r^{n+1}-1}{r-1}
\]

\end{document}
